% notes/artefacts/execution-market-making/order-types-instructions.tex
% Legacy file kept verbatim when it was migrated (notes/source-review-records.md, order matching and fragmentation / source review / optimal execution / execution control reviews): the whole of
% compendium/legacy/execution-market-making/02_orders_matching_and_limit_order_books/01_order_types_and_instructions.tex
% below its migrate-to line (migrate-to: execution-market-making, market-making-lob, market-structure-execution-quality), with only the subsection heading unnumbered. The compendium entries
% state each claim with its locator; the original wording lives here. Treat everything here as unchecked own working.
% Destinations: Order, Order-Instruction Segmentation, Core Order-Type Segmentation, Market Order Cost Segmentation, Conditional Activation, Stops Are Not Limits, Tick-Sensitive Orders, Time-in-Force, Quantity and Handling, Special Instruction, Instruction Semantics -> def:order-instructions; Marketable Limit Order -> def:limit-order-book and def:order-instructions; Order Choice Trades Certainty for Leakage, Market vs Limit Order Choice, Limit-Order Fill Threshold (with its proof), Spread-Only Rules Fail -> res:limit-order-options; from Harris chap. 4 (pp. 68-86), chap. 6 (pp. 117, 126), chap. 14 (p. 304), Bouchaud et al. p. 55, FIX 4.4 code sets.
% Compile inside notes/artefacts/execution-market-making/ with  latexmk -pdf order-types-instructions.tex
\PassOptionsToPackage{table}{xcolor}
\documentclass[a4paper,oneside,openany]{book}
% Your own packages, macros and environments. The learning guide's entry environments,
% \citep and \qfsource are in house.tex, which compendium.tex loads after this file.

% -----------------------------
% Document Class & Encoding
% -----------------------------
\usepackage[utf8]{inputenc}  % Input encoding
\usepackage[T1]{fontenc}     % Font encoding

% -----------------------------
% Page Layout & Geometry
% -----------------------------
\usepackage[margin=2cm,foot=1cm]{geometry}

% -----------------------------
% Header & Footer Management
% -----------------------------
\usepackage{fancyhdr}
\pagestyle{fancy}

% -----------------------------
% Text Formatting & Spacing
% -----------------------------
\usepackage{ragged2e}
\usepackage[nodisplayskipstretch]{setspace}
\setstretch{1}  % Line spacing

% -----------------------------
% Listings for Code
% -----------------------------
\usepackage{listings}
\usepackage{xcolor}

% -----------------------------
% Tables, Figures & Graphics
% -----------------------------
\usepackage{tabularx}
\usepackage{multirow}
\usepackage{pdflscape}
\usepackage{graphicx}

% -----------------------------
% Algorithms
% -----------------------------
\usepackage{algorithm}       % For floating algorithm environment
\usepackage{algpseudocode}   % For pseudo-code inside algorithms

% -----------------------------
% Math Packages
% -----------------------------
\usepackage{amsmath,amssymb,amsthm,mathtools,bm}
\usepackage{mathrsfs}        % Script fonts
\usepackage{tikz-cd}         % Commutative diagrams
\usepackage{rotating}        % Rotated objects
\usepackage{float}           % Float placement
\usetikzlibrary{calc, arrows, positioning, arrows.meta}

% -----------------------------
% TikZ and its Libraries
% -----------------------------
\usepackage{tikz}
\usetikzlibrary{positioning} % Added to support "below=of" syntax
\usepackage{forest}          % For advanced tree diagrams
\usepackage{tikz-qtree}      % For vertical tree diagrams

% -----------------------------
% Miscellaneous Packages
% -----------------------------
\usepackage{xparse}          % For defining robust commands
\usepackage{enumitem}        % For customizing list layouts
\usepackage{makecell}        % For enhanced table cells
\usepackage{comment}		 % For commenting out sections

%%%%%%%%%%%%%%%%%%%%%%%%%%%%%%%%%%%%%%%%%%%%%%%%%%%%%%%%%%%%%%%%%%%%%%%%%%%%%%%
% Custom Math Commands and Shortcuts
%%%%%%%%%%%%%%%%%%%%%%%%%%%%%%%%%%%%%%%%%%%%%%%%%%%%%%%%%%%%%%%%%%%%%%%%%%%%%%%

% Matrix and Vector Formatting
\newcommand{\matr}[1]{\mathbf{#1}}      % For undergraduate algebra (alternative: \bm{#1})
\newcommand{\vect}[1]{\bm{#1}}

% Standard Number Sets
\newcommand{\N}{\mathbb{N}}
\newcommand{\R}{\mathbb{R}}
\newcommand{\Q}{\mathbb{Q}}
\newcommand{\C}{\mathbb{C}}
\newcommand{\Z}{\mathbb{Z}}

% Bold variables (customize as needed)
\newcommand{\x}{\mathbf{x}}
\newcommand{\F}{\mathbf{F}}
\newcommand{\f}{\mathbf{f}}
\newcommand{\y}{\mathbf{y}}
\renewcommand{\a}{\mathbf{a}}
\renewcommand{\b}{\mathbf{b}}
\renewcommand{\c}{\mathbf{c}}
\newcommand{\h}{\mathbf{h}}
\newcommand{\g}{\mathbf{g}}
\newcommand{\z}{\mathbf{z}}
\newcommand{\ze}{\mathbf{0}}

% Norms, Absolute Values, and Bracketing
\newcommand{\norm}[1]{\left\lVert#1\right\rVert}
\newcommand{\abs}[1]{\left\lvert#1\right\rvert}
\newcommand{\brk}[1]{\left[#1\right]}
\newcommand{\brc}[1]{\left\{#1\right\}}
\newcommand{\paren}[1]{\left(#1\right)}
\newcommand{\normop}[1]{\left\lVert#1\right\rVert_\text{op}}

% Miscellaneous Math Shortcuts
\newcommand{\LL}{\mathcal{L}}
\newcommand{\uni}{\overset{\text{uni}}{\to}}
\DeclareMathOperator{\diam}{diam}
\newcommand{\Prr}[1]{\text{Pr}\left(#1\right)}
\newcommand{\ceil}[1]{\lceil #1 \rceil}
\newcommand{\floor}[1]{\lfloor #1 \rfloor}
\newcommand\disteq{\stackrel{\textnormal{dist}}{=}}

% Evaluation command with options
\NewDocumentCommand{\evalat}{sO{\big}mm}{%
  \IfBooleanTF{#1}
    {\mleft. #3 \mright|_{#4}}
    {#3#2|_{#4}}%
}

%%%%%%%%%%%%%%%%%%%%%%%%%%%%%%%%%%%%%%%%%%%%%%%%%%%%%%%%%%%%%%%%%%%%%%%%%%%%%%%
% Color Definitions
%%%%%%%%%%%%%%%%%%%%%%%%%%%%%%%%%%%%%%%%%%%%%%%%%%%%%%%%%%%%%%%%%%%%%%%%%%%%%%%

\definecolor{dartmouthgreen}{rgb}{0.05, 0.5, 0.06}
\definecolor{egyptianblue}{rgb}{0.06, 0.2, 0.65}
\definecolor{dukeblue}{rgb}{0.0, 0.0, 0.61}
\definecolor{jazzberryjam}{rgb}{0.65, 0.04, 0.37}
\definecolor{magenta}{HTML}{EC008C}
\definecolor{darkmagenta}{rgb}{0.55, 0.0, 0.55}
\definecolor{deeppink}{rgb}{1.0, 0.08, 0.58}
\definecolor{codegreen}{rgb}{0,0.6,0}
\definecolor{codegray}{rgb}{0.5,0.5,0.5}
\definecolor{codepurple}{rgb}{0.58,0,0.82}
\definecolor{backcolour}{rgb}{0.95,0.95,0.92}

%%%%%%%%%%%%%%%%%%%%%%%%%%%%%%%%%%%%%%%%%%%%%%%%%%%%%%%%%%%%%%%%%%%%%%%%%%%%%%%
% Custom Shortcuts and Emphasis Commands
%%%%%%%%%%%%%%%%%%%%%%%%%%%%%%%%%%%%%%%%%%%%%%%%%%%%%%%%%%%%%%%%%%%%%%%%%%%%%%%

\newcommand{\isomorp}{\xrightarrow{\sim}}
\newcommand{\hlt}[1]{\textit{{\color{blue}#1}}}
\newcommand{\simpt}[1]{\textit{{\color{deeppink}#1}}}
\newcommand{\impt}[1]{\textit{{\color{red}#1}}}
\newcommand{\gcds}[1]{\textnormal{gcd}#1}
\newcommand{\mins}[1]{\textnormal{min}#1}
\newcommand{\maxs}[1]{\textnormal{max}#1}
\newcommand{\lcms}[1]{\textnormal{lcm}#1}
\newcommand{\degs}[1]{\textnormal{deg}#1}
\newcommand{\exps}[1]{\textnormal{exp}#1}
\newcommand{\tors}[1]{\textnormal{Tor}#1}
\newcommand{\homs}[1]{\textnormal{Hom}#1}
\newcommand{\anns}[1]{\textnormal{Ann}#1}
\newcommand{\defeq}{\vcentcolon=}
\DeclareMathOperator{\sign}{sign}

%%%%%%%%%%%%%%%%%%%%%%%%%%%%%%%%%%%%%%%%%%%%%%%%%%%%%%%%%%%%%%%%%%%%%%%%%%%%%%%
% Theorem and Definition Environments
%%%%%%%%%%%%%%%%%%%%%%%%%%%%%%%%%%%%%%%%%%%%%%%%%%%%%%%%%%%%%%%%%%%%%%%%%%%%%%%
\theoremstyle{plain}
\newtheorem{theorem}{Theorem}[section]
\newtheorem{corollary}[theorem]{Corollary}
\newtheorem{lemma}[theorem]{Lemma}
\newtheorem{proposition}[theorem]{Proposition}
\newtheorem{qnbank}[theorem]{Question Bank}
\let\oldqnbank\qnbank
\renewcommand{\qnbank}{\oldqnbank\normalfont}
\newtheorem{process}[theorem]{Process}
\let\oldprocess\process
\renewcommand{\process}{\oldprocess\normalfont}
\newtheorem{method}[theorem]{Method}
\let\oldmethod\method
\renewcommand{\method}{\oldmethod\normalfont}
\newtheorem{definition}[theorem]{Definition}
\let\olddefinition\definition
\renewcommand{\definition}{\olddefinition\normalfont}
\newtheorem{example}[theorem]{Example}
\let\oldexample\example
\renewcommand{\example}{\oldexample\normalfont}
\newtheorem{remark}[theorem]{Remark}
\let\oldremark\remark
\renewcommand{\remark}{\oldremark\normalfont}
\newtheorem{notation}[theorem]{Notation}
\let\oldnotation\notation
\renewcommand{\notation}{\oldnotation\normalfont}
\newtheorem{principle}[theorem]{Principle}
\let\oldprinciple\principle
\renewcommand{\principle}{\oldprinciple\normalfont}
\newtheorem{mechanism}[theorem]{Mechanism}
\let\oldmechanism\mechanism
\renewcommand{\mechanism}{\oldmechanism\normalfont}
\newtheorem{taxonomy}[theorem]{Taxonomy}
\let\oldtaxonomy\taxonomy
\renewcommand{\taxonomy}{\oldtaxonomy\normalfont}
\newtheorem{model}[theorem]{Model}
\let\oldmodel\model
\renewcommand{\model}{\oldmodel\normalfont}
\newtheorem{assumption}[theorem]{Assumption}
\let\oldassumption\assumption
\renewcommand{\assumption}{\oldassumption\normalfont}
\newtheorem{casestudy}[theorem]{Case Study}
\let\oldcasestudy\casestudy
\renewcommand{\casestudy}{\oldcasestudy\normalfont}
\newtheorem{empiricalfinding}[theorem]{Empirical Finding}
\let\oldempiricalfinding\empiricalfinding
\renewcommand{\empiricalfinding}{\oldempiricalfinding\normalfont}
\newtheorem{practitionernote}[theorem]{Practitioner Note}
\let\oldpractitionernote\practitionernote
\renewcommand{\practitionernote}{\oldpractitionernote\normalfont}
\newtheorem{caveat}[theorem]{Caveat}
\let\oldcaveat\caveat
\renewcommand{\caveat}{\oldcaveat\normalfont}
\newtheorem{marketrule}[theorem]{Rule}
\let\oldmarketrule\marketrule
\renewcommand{\marketrule}{\oldmarketrule\normalfont}
\newtheorem{takeaway}[theorem]{Takeaway}
\let\oldtakeaway\takeaway
\renewcommand{\takeaway}{\oldtakeaway\normalfont}

%%%%%%%%%%%%%%%%%%%%%%%%%%%%%%%%%%%%%%%%%%%%%%%%%%%%%%%%%%%%%%%%%%%%%%%%%%%%%%%
% Breakable Algorithm Environment (Algorithms spanning multiple pages)
%%%%%%%%%%%%%%%%%%%%%%%%%%%%%%%%%%%%%%%%%%%%%%%%%%%%%%%%%%%%%%%%%%%%%%%%%%%%%%%

\makeatletter
\newenvironment{breakablealgorithm}
 {%
  \begin{center}
    \refstepcounter{algorithm} % Increment algorithm counter
    \hrule height.8pt depth0pt \kern2pt
    \renewcommand{\caption}[2][\relax]{%
      {\raggedright\textbf{Algorithm \thealgorithm:} ##2\par}%
      \ifx\relax##1\relax
        \addcontentsline{loa}{algorithm}{\protect\numberline{\thealgorithm}##2}%
      \else
        \addcontentsline{loa}{algorithm}{\protect\numberline{\thealgorithm}##1}%
      \fi
      \kern2pt\hrule\kern2pt
    }%
 }{%
    \kern2pt\hrule
  \end{center}
 }
\makeatother

%%%%%%%%%%%%%%%%%%%%%%%%%%%%%%%%%%%%%%%%%%%%%%%%%%%%%%%%%%%%%%%%%%%%%%%%%%%%%%%
% Code Listings Style
%%%%%%%%%%%%%%%%%%%%%%%%%%%%%%%%%%%%%%%%%%%%%%%%%%%%%%%%%%%%%%%%%%%%%%%%%%%%%%%

\lstdefinestyle{codestyle}{
  commentstyle=\color{codegreen},
  keywordstyle=\color{magenta},
  numberstyle=\tiny\color{codegray},
  stringstyle=\color{codepurple},
  basicstyle=\ttfamily\normalsize,
  breakatwhitespace=false,
  breaklines=true,
  captionpos=b,
  keepspaces=true,
  numbers=left,
  numbersep=10pt,
  frame=single,
  showspaces=false,
  showstringspaces=false,
  showtabs=false,
  tabsize=2
}
\lstset{style=codestyle}

% -----------------------------
% Hyperlinks (loaded last: loaded before float and algorithm it duplicated figure anchors)
% -----------------------------
\usepackage{hyperref}
\hypersetup{
  colorlinks,
  linkcolor={blue!80!black},
  citecolor={green!80!black},
  urlcolor={red!80!black}
}

\usepackage[T1]{fontenc}
\usepackage{amsmath,amssymb,booktabs,graphicx,hyperref}
\hypersetup{colorlinks}
\setlength{\emergencystretch}{3em}   % lets a line holding a long key break
\newcounter{qfentry}[chapter]
\renewcommand{\theqfentry}{\thechapter.\arabic{qfentry}}
% \qftype{type}{Name} declares the environment qf<type>; its title wraps like a paragraph.
\newcommand{\qftype}[2]{\newenvironment{qf#1}[1][]{\par\addvspace{\medskipamount}%
  \refstepcounter{qfentry}{\raggedright\noindent\bfseries #2 \theqfentry
  \if\relax\detokenize{##1}\relax\else\quad ##1\fi\par}\nopagebreak}%
  {\par\addvspace{\medskipamount}}}
\qftype{definition}{Definition}\qftype{convention}{Convention}\qftype{result}{Result}
\qftype{model}{Model}\qftype{method}{Method}\qftype{mechanism}{Mechanism}
\qftype{evidence}{Evidence}\qftype{data}{Data}\qftype{rule}{Rule}
\qftype{pitfall}{Pitfall}\qftype{example}{Example}
\newcommand{\qffield}[1]{\par\textbf{#1.}\enspace}   % starts one of an entry's fields
% \citep[locator]{key} links the key to its line in the sources if it has one,
% and otherwise to its entry in the pinned catalogue PDF.
\newcommand{\qfcatalogue}{../guide/institutional_quantitative_finance_reading_list.pdf}
\NewDocumentCommand{\citep}{o o m}{[\IfValueT{#2}{#1 }%
  \ifcsname r@src:#3\endcsname\hyperref[src:#3]{\nolinkurl{#3}}%
  \else\href{\qfcatalogue\##3}{\nolinkurl{#3}}\fi
  \IfValueTF{#2}{, #2}{\IfValueT{#1}{, #1}}]}
% \qfsource{key}{reference} records one source outside the catalogue, in front/sources.tex.
\newcommand{\qfsource}[2]{\par\noindent\phantomsection\label{src:#1}\nolinkurl{#1}\quad #2\par}
% \qfsections{key, ...} inputs <key>/section.tex for each listed key that has one, in order.
\ExplSyntaxOn
\NewDocumentCommand{\qfsections}{m}
  {\clist_map_inline:nn {#1} {\InputIfFileExists{##1/section}{}{}}}
\ExplSyntaxOff

\pagestyle{plain}
\begin{document}
\chapter*{Legacy notes: order types and instructions}

% From compendium/legacy/execution-market-making/02_orders_matching_and_limit_order_books/01_order_types_and_instructions.tex (whole file, deleted 2026-10-16 once empty).
\subsection*{Order Types and Instructions}

\begin{definition} \hlt{Order}\\
An order is a trade instruction. At minimum it identifies the instrument, side, and quantity; practical orders also specify price limits, activation rules, validity, display, routing, settlement, or broker discretion.
\end{definition}

\begin{taxonomy} \hlt{Order-Instruction Segmentation}
\begin{enumerate}[label=\roman*.]
\setlength{\itemsep}{0pt}
\item \hlt{economic terms}: instrument, side, quantity, and price constraint;
\item \hlt{timing}: submission time, activation trigger, validity, and expiry;
\item \hlt{execution constraint}: immediate execution, full execution, minimum size, or partial-fill permission;
\item \hlt{display}: fully displayed, hidden, reserve, or iceberg quantity;
\item \hlt{handling}: held by strict rules, or not-held with broker discretion;
\item \hlt{post-trade terms}: settlement, clearing, and account capacity.
\end{enumerate}
The label of an order is therefore only a shorthand for a bundle of instructions.
\end{taxonomy}

\begin{taxonomy} \hlt{Core Order-Type Segmentation}
\begin{enumerate}[label=\roman*.]
\setlength{\itemsep}{0pt}
\item \hlt{market}: pay immediacy and consume the best available liquidity;
\item \hlt{limit}: specify a worst acceptable price and wait in the book if not marketable;
\item \hlt{stop}: become active only after a trigger price is reached;
\item \hlt{market-if-touched}: become active when price moves favourably to the touch price;
\item \hlt{IOC/FOK}: cancel any unfilled quantity immediately, or require full immediate execution;
\item \hlt{pegged}: reference the midpoint, best bid, best offer, or another venue-derived price;
\item \hlt{hidden/iceberg}: reduce displayed size while keeping some executable interest undisplayed.
\end{enumerate}
\end{taxonomy}

\begin{principle} \hlt{Order Choice Trades Certainty for Leakage}\\
Aggressive orders buy execution certainty but reveal urgency. Passive and hidden orders reduce signalling, but they introduce queue risk, missed fills, and adverse selection.
\end{principle}

\begin{mechanism} \hlt{Market Order Cost Segmentation}
\begin{enumerate}[label=\roman*.]
\setlength{\itemsep}{0pt}
\item \hlt{spread cost}: a small market buy normally pays the ask and a small market sell receives the bid;
\item \hlt{price improvement}: negotiation, midpoint interaction, or discretionary handling may beat displayed quote;
\item \hlt{market impact}: large orders walk the book and move prices to attract the other side;
\item \hlt{execution-price uncertainty}: the quote can change between submission and execution.
\end{enumerate}
The market order solves execution uncertainty by accepting price uncertainty.
\end{mechanism}

\begin{definition} \hlt{Marketable Limit Order}\\
A marketable limit order is priced aggressively enough to trade immediately if available liquidity is still present. It is like a market order with a price cap for buys or a price floor for sells; unfilled residual quantity either rests or cancels according to its other instructions.
\end{definition}

\begin{principle} \hlt{Market vs Limit Order Choice}\\
Market orders exchange price control for execution certainty. Limit orders exchange execution certainty for price control. Impatient traders, informed traders with decaying signals, and traders facing narrow spreads tend to cross the spread; patient traders, liquidity suppliers, and traders facing wide spreads tend to post.
\end{principle}

\begin{proposition} \hlt{Limit-Order Fill Threshold}\\
For a buy order, let $p$ be the probability that a passive limit order fills, $S$ the spread saved if it fills, $M$ the expected cost of chasing the trade if it does not fill, and $A$ the adverse-selection cost conditional on being filled. Posting is preferred to crossing when $p(S-A)>(1-p)M$.
\begin{proof}
Relative to crossing immediately, posting gains $S$ and loses $A$ when it fills.\\
If it misses, the trader pays expected chase cost $M$. The expected relative gain is $p(S-A)-(1-p)M$, so posting is preferred exactly when this is positive.
\end{proof}
\end{proposition}

\begin{caveat} \hlt{Spread-Only Rules Fail}\\
Rules such as ``cross when spreads are narrow, post when spreads are wide'' work only when the trader has no valuation information. A trader who knows the current quote is stale may rationally demand liquidity even when the spread looks wide, or supply liquidity even when the spread looks narrow.
\end{caveat}

\begin{taxonomy} \hlt{Conditional Activation Segmentation}
\begin{enumerate}[label=\roman*.]
\setlength{\itemsep}{0pt}
\item \hlt{stop order}: activates after adverse price move; buy stops trigger after price rises, sell stops after price falls;
\item \hlt{stop-limit order}: activates after the stop, but can execute only within the specified limit price;
\item \hlt{market-if-touched order}: activates after a favourable price move, then behaves like a market order.
\end{enumerate}
Stops tend to demand liquidity in the direction of the move; market-if-touched orders are more contrarian.
\end{taxonomy}

\begin{caveat} \hlt{Stops Are Not Limits}\\
A stop price controls activation, not execution price. A stop-limit adds a limit price after activation, but then loses execution certainty when the market gaps through the limit.
\end{caveat}

\begin{marketrule} \hlt{Tick-Sensitive Orders}\\
Tick-sensitive orders condition execution on the direction of the last price change. A buy downtick order waits for a downtick or zero downtick; a sell uptick order waits for an uptick or zero uptick. These instructions are mainly useful when tick sizes are large enough that one tick is economically meaningful.
\end{marketrule}

\begin{taxonomy} \hlt{Time-in-Force Segmentation}
\begin{enumerate}[label=\roman*.]
\setlength{\itemsep}{0pt}
\item \hlt{day}: valid until the session close;
\item \hlt{good-till-cancel}: valid until the trader or broker cancels it;
\item \hlt{good-till-date}: valid until a specified date or time;
\item \hlt{IOC}: execute immediately and cancel the unfilled balance;
\item \hlt{FOK}: execute the full quantity immediately or cancel;
\item \hlt{market-on-open/close}: execute in the opening or closing auction.
\end{enumerate}
Validity instructions decide when an order can become stale enough to disappear.
\end{taxonomy}

\begin{taxonomy} \hlt{Quantity and Handling Segmentation}
\begin{enumerate}[label=\roman*.]
\setlength{\itemsep}{0pt}
\item \hlt{all-or-none}: the whole quantity must execute;
\item \hlt{minimum acceptable quantity}: each fill must meet a specified minimum size;
\item \hlt{display size}: only part of the order is shown;
\item \hlt{not-held}: the broker may use discretion over timing, price, and exposure.
\end{enumerate}
These instructions trade fill probability against information leakage and accountability.
\end{taxonomy}

\begin{taxonomy} \hlt{Special Instruction Segmentation}
\begin{enumerate}[label=\roman*.]
\setlength{\itemsep}{0pt}
\item \hlt{spread orders}: trade two related instruments as one package, usually with a limit on the price difference;
\item \hlt{substitution orders}: let the broker choose among acceptable substitutes;
\item \hlt{special settlement}: request nonstandard settlement, often around record dates or corporate actions.
\end{enumerate}
\end{taxonomy}

\begin{caveat} \hlt{Instruction Semantics Are Venue-Specific}\\
The same label can behave differently across venues. Minimum quantity, display refresh, priority loss after modification, pegging reference, and cancellation rules must be read as part of the market design.
\end{caveat}

\end{document}
